%% sections/introduction.tex

\section{Introduction}
\label{sec:introduction}

The rapid proliferation of large language models (LLMs) has fundamentally transformed the landscape of text generation. Systems such as GPT-4, Claude, and Gemini can now produce fluent, coherent, and contextually rich prose that is, in many cases, indistinguishable from human writing. While these capabilities carry substantial scientific and commercial value, they also introduce serious challenges for content authenticity, academic integrity, and information reliability. As LLMs become increasingly accessible, the ability to automatically detect machine-generated text has emerged as a critical research priority.

The problem is far from solved. Early detection systems exploited the statistical regularities left by language models—unnaturally low perplexity, characteristic token distributions, and uniform sentence rhythms—to separate human from machine text with reasonable accuracy. However, adversarial text transformations, including synonym substitution, paraphrasing, and character-level obfuscation, have shown that surface-level detectors are brittle and easily circumvented~\cite{sadasivan:2023}. When LLMs are additionally instructed to mimic the style of a specific human author, detection becomes considerably harder, since the model no longer relies on its default generation patterns~\cite{jakesch:2023}.

This setting—adversarially styled, obfuscated AI-generated text—is precisely the focus of the \textit{Voight-Kampff Generative AI Detection} shared task at PAN 2026. Building on two prior editions of the challenge~\cite{bevendorff:2024d,bevendorff:2025d}, participants receive text samples that may have been produced by a human or by an LLM instructed to imitate human style, and must produce a confidence score indicating the likelihood of machine authorship. The test set deliberately includes novel models and unknown obfuscation strategies to probe robustness.

Against this backdrop, we present \textbf{RADAR} (\textbf{R}obust \textbf{A}dversarial-Resistant \textbf{A}I Text \textbf{D}etection \textbf{A}rchitecture), a multimodal detection system that draws on two complementary sources of evidence. The first is \textit{semantic representation}: RADAR uses ModernBERT-large~\cite{modernbert:2024}—a state-of-the-art bidirectional encoder pre-trained on a large and diverse corpus—to capture deep contextual meaning from the input text. The second is \textit{stylometric analysis}: a custom Statistical Linguistic Profile Extractor (\slpe) computes 38 handcrafted features spanning vocabulary richness, syntactic complexity, discourse coherence, style fingerprinting, and word-level entropy. These two streams are fused through a cross-attention mechanism, and the resulting representation is passed to a binary classification head.

The key intuition behind RADAR is that adversarial obfuscations can substantially alter the surface form of text—and thereby defeat detectors that rely purely on token statistics or LM perplexity—yet they leave systematic traces at the structural, discourse, and stylistic levels that are difficult to simultaneously hide. By fusing semantic and stylometric evidence, RADAR targets a richer feature space that is inherently more resistant to single-mode attacks.

We train RADAR using a two-stage curriculum: in the first stage, the encoder weights are frozen and only the stylometric projector, fusion layer, and classifier are trained; in the second stage, all weights are jointly fine-tuned end-to-end. This approach avoids catastrophic forgetting of the encoder's pre-trained representations while allowing the model to adapt its semantic features to the detection task.

Our experiments demonstrate that RADAR achieves strong performance on both the official PAN-2026 validation set and a challenging out-of-domain benchmark assembled from the RAID corpus~\cite{dugan:2024}. The final end-to-end model reaches a mean score of 0.978 on the combined validation set and 0.972 on the RAID out-of-domain test set, substantially outperforming all provided baselines.

The remainder of this paper is organised as follows. Section~\ref{sec:relatedwork} reviews related work on AI-generated text detection. Section~\ref{sec:task} describes the PAN-2026 task and its evaluation protocol. Section~\ref{sec:methodology} presents the RADAR architecture in detail. Section~\ref{sec:experiments} describes our experimental setup, datasets, and training procedure. Section~\ref{sec:results} reports and analyses the results. Section~\ref{sec:conclusion} concludes.
